%!TEX program = xelatex
\let\stop\empty

\documentclass[12pt,a4paper]{exam-zh}

\usepackage{xeCJK}
\setCJKmainfont{Noto Serif CJK SC}
\usepackage{geometry}
\usepackage{array}
\usepackage{tikz}
\usepackage{amsmath}
\usepackage{wrapfig}
\usepackage{adjustbox}
\usepackage{enumitem}
\usepackage{fancyhdr}
\usepackage{pgfplots}
\pgfplotsset{compat=1.18}
\usetikzlibrary{calc}

\usepackage{scalerel}
\def\mcirc{\mathbin{\scalerel*{\circ}{j}}}
\def\msquare{\mathord{\scalerel*{\Box}{gX}}}
\geometry{margin=2cm}

\AtBeginDocument{\zihao{4}}
% ----------- 调整 exam -----------


% --- 全局样式定义 ---
% 定义选择题环境：题号 1. 2.，选项 A. B. C. D. 粗体，且保留 1.5em 间距
\examsetup{
    page = {
    size = a4paper,           % 设定纸张大小 [4]
    foot-content = { - ; - }  % 页脚显示为 -1-, -2- 格式 [3]
  },
    choices = {
    label-pos = top-left,
    column-sep = 1em,
    label-pos = auto,
    label-sep = 0.5em,
    max-columns = 4
    },
  question = {
    top-sep    = 1em, % 调整题目上方的垂直距离（默认值为 .5em plus .5em minus .2em）
    bottom-sep = 1em  % 调整题目下方的垂直距离（默认值为 .5em plus .5em minus .2em）
  },
    paren = {
    show-paren = false  % 设置为 false 则不显示括号 [3, 4]
  },
    choices = {
    label = \textbf{\Alph*}
  }, 
    solution = {
    show-solution = hide,      % 不生成空答案页
    label-content = {参考答案},  % 答案区的标题文字
    show-qed = true            % 在解答末尾显示结束符号 ◻ 
  }
}

\newcommand{\functionblock}[6]{%
  \node[anchor=north west, inner sep=0pt] at
    ($(current page.north west)+(#1\paperwidth/2,-#2\paperheight/6)$) {%
      \begin{minipage}[c][\paperheight/6][c]{\paperwidth/2}
        \centering
        {\large $#3$}\par\vspace{0.1em}
        \begin{tikzpicture}
          \begin{axis}[
            width=0.34\paperwidth,
            height=0.085\paperheight,
            scale only axis,
            axis lines=middle,
            axis line style={->},
            xmin=#4, xmax=#5,
            ymin=#4, ymax=#5,
            enlargelimits=false,
            xtick=\empty,
            ytick=\empty,
            tick style={draw=none},
            samples=160,
            unbounded coords=jump,
            clip=true
          ]
          \end{axis}
        \end{tikzpicture}
      \end{minipage}%
    };%
}


\newcommand{\fourfunctionblock}[6]{%
  \node[anchor=north west, inner sep=0pt] at
    ($(current page.north west)+(#1\paperwidth/2,-#2\paperheight/2)$) {%
      \begin{minipage}[c][\paperheight/2][c]{\paperwidth/2}
        \centering
        {\Large $#3$}\par\vspace{0.8em}
        \begin{tikzpicture}
          \begin{axis}[
            width=0.34\paperwidth,
            height=0.24\paperheight,
            scale only axis,
            axis lines=middle,
            axis line style={->},
            xmin=#4, xmax=#5,
            ymin=#4, ymax=#5,
            enlargelimits=false,
            xtick=\empty,
            ytick=\empty,
            tick style={draw=none},
            samples=160,
            unbounded coords=jump,
            clip=true
          ]
          \end{axis}
        \end{tikzpicture}
      \end{minipage}%
    };%
}

\begin{document}

\begin{tikzpicture}[remember picture, overlay]
  \foreach \i in {1,...,5} {
    \draw ($(current page.north west)!\i/6!(current page.south west)$) --
          ($(current page.north east)!\i/6!(current page.south east)$);
  }
  \draw ($(current page.north west)!1/2!(current page.north east)$) --
        ($(current page.south west)!1/2!(current page.south east)$);

  \functionblock{0}{0}{y=x}{-3}{3}{\addplot[thick, blue, domain=-3:3] {x};}
  \functionblock{1}{0}{y=x^2}{-3}{3}{\addplot[thick, blue, domain=-3:3] {x^2};}
  \functionblock{0}{1}{y=x^3}{-3}{3}{\addplot[thick, blue, domain=-3:3] {x^3};}
  \functionblock{1}{1}{y=\dfrac{1}{x}}{-3}{3}{\addplot[thick, blue, domain=-3:-0.15] {1/x}; \addplot[thick, blue, domain=0.15:3] {1/x};}
  \functionblock{0}{2}{y=\dfrac{1}{x^2}}{-3}{3}{\addplot[thick, blue, domain=-3:-0.15] {1/(x^2)}; \addplot[thick, blue, domain=0.15:3] {1/(x^2)};}
  \functionblock{1}{2}{y=\sqrt{x}}{-0.5}{4}{\addplot[thick, blue, domain=0:4] {sqrt(x)};}
  \functionblock{0}{3}{y=-x}{-3}{3}{\addplot[thick, red, domain=-3:3] {-x};}
  \functionblock{1}{3}{y=-x^2}{-3}{3}{\addplot[thick, red, domain=-3:3] {-x^2};}
  \functionblock{0}{4}{y=-x^3}{-3}{3}{\addplot[thick, red, domain=-3:3] {-x^3};}
  \functionblock{1}{4}{y=-\dfrac{1}{x}}{-3}{3}{\addplot[thick, red, domain=-3:-0.15] {-1/x}; \addplot[thick, red, domain=0.15:3] {-1/x};}
  \functionblock{0}{5}{y=-\dfrac{1}{x^2}}{-3}{3}{\addplot[thick, red, domain=-3:-0.15] {-1/(x^2)}; \addplot[thick, red, domain=0.15:3] {-1/(x^2)};}
  \node[anchor=north west, inner sep=0pt] at
    ($(current page.north west)+(1\paperwidth/2,-5\paperheight/6)$) {%
      \begin{minipage}[c][\paperheight/6][c]{\paperwidth/2}
        \centering
        {\large $y=-\sqrt{x}$}\par\vspace{0.1em}
        \begin{tikzpicture}
          \begin{axis}[
            width=0.34\paperwidth,
            height=0.085\paperheight,
            scale only axis,
            axis lines=middle,
            axis line style={->},
            xmin=-0.5, xmax=4,
            ymin=-3, ymax=1,
            enlargelimits=false,
            xtick=\empty,
            ytick=\empty,
            tick style={draw=none},
            samples=160,
            unbounded coords=jump,
            clip=true
          ]
          \end{axis}
        \end{tikzpicture}
      \end{minipage}%
    };
\end{tikzpicture}


\newpage

\begin{tikzpicture}[remember picture, overlay]
  \draw ($(current page.north west)!1/2!(current page.south west)$) --
        ($(current page.north east)!1/2!(current page.south east)$);
  \draw ($(current page.north west)!1/2!(current page.north east)$) --
        ($(current page.south west)!1/2!(current page.south east)$);

  \fourfunctionblock{0}{0}{y=\dfrac{1}{(x+2)^2}}{-6}{6}{}
  \fourfunctionblock{1}{0}{y=-\dfrac{1}{2(x-1)^2}}{-6}{6}{}
  \fourfunctionblock{0}{1}{y=3\sqrt{x+2}-4}{-6}{6}{}
  \fourfunctionblock{1}{1}{y=\dfrac{4}{2x+3}}{-6}{6}{}
\end{tikzpicture}

\end{document}